% Requires: \usepackage{booktabs}
% Fill \cite{} placeholders. Table~\ref{tab:sparseloop} rows marked --- are pending
% the overnight Sparseloop sweep (sparse_sweep_results.txt).

\subsection{From Reported Sparsity to a Realized Speed-Up}
\label{subsec:sparse-harvest}

The pruned networks of the preceding section report weight sparsities of
$85$--$99\%$, yet none of them run any faster than their dense counterparts on
commodity hardware. The reason is that magnitude pruning produces
\emph{unstructured} sparsity: the surviving weights occupy irregular positions,
so a general-purpose matrix multiply---on a CPU, a GPU, or through an inference
runtime such as ONNX---still streams and multiplies the zeros. The floating-point
operation count is unchanged, and so are latency and energy. Turning a reported
sparsity into a measurable benefit therefore requires either (i)~restructuring the
sparsity so that dense hardware can skip it, or (ii)~a specialized accelerator that
skips zeros natively. We evaluate both.

\paragraph{Structured channel harvesting.}
We observe that at sufficiently high sparsity, unstructured magnitude pruning
\emph{emergently} drives entire output channels to zero: for the $98\%$-sparse
VGG-19, $70\%$ of all convolutional filters are identically zero. Such channels can
be deleted outright---together with the corresponding input channels they feed---to
yield a smaller \emph{dense} network that runs faster on any hardware, with no
special kernels. Because the removed channels contribute nothing to the output
(after folding the batch-normalization/bias constants they would otherwise emit),
the transformation is \emph{exactly lossless}: the compact model is numerically
identical to the original (maximum output deviation $<10^{-12}$ in double
precision), so accuracy is preserved with zero fine-tuning. Table~\ref{tab:harvest}
reports the result. The high-sparsity feed-forward networks collapse by
$3.9$--$12.1\times$ in parameters and deliver up to $4.2\times$ measured GPU
throughput, entirely losslessly.

Two limits of the method are equally important. First, the effect is subject to a
\emph{sparsity threshold} near $95$--$97\%$: below it, unstructured pruning does not
kill whole channels (VGG-19 at $90\%$ sparsity harvests to $1.0\times$), so the
compression is negligible. Second, and more fundamentally, the method is
\emph{architecture dependent}. Residual and densely-connected networks distribute
sparsity across their skip and concatenation paths, so individual channels almost
never die completely; a dependency-aware pruning of every ResNet-32 and DenseNet-40
checkpoint removes \emph{zero} channels losslessly, leaving them at $1.0\times$
(Table~\ref{tab:harvest}, lower block). For these architectures, structured
harvesting offers nothing.

\paragraph{Accelerator-level exploitation.}
For the sparsity that structured harvesting cannot reach---the residual networks in
full, and the $\sim\!87\%$ unstructured sparsity that remains \emph{inside} the kept
channels of the harvested feed-forward networks---we model a sparsity-aware
accelerator with the Sparseloop~\cite{sparseloop}/Timeloop~\cite{timeloop} analytical
framework: a $256$-PE weight-stationary array whose per-processing-element datapath
\emph{skips} multiply--accumulate operations on zero weights.
Table~\ref{tab:sparseloop} summarizes the outcome. On the $98\%$-sparse VGG-19 the
accelerator converts the residual sparsity into a $3.68\times$ whole-network
throughput improvement and a $21\%$ energy reduction---gains that a dense
processor structurally cannot obtain. The result carries two caveats that shape the
broader picture. The throughput benefit only materializes when the mapper is
permitted to trade energy for latency and when zero-\emph{skipping} (rather than
mere power-\emph{gating}) is modeled; and weight compression is only admissible at
the innermost, per-PE register, because compressing weights above the spatial array
violates the array's fan-out and renders every dataflow mapping invalid. In other
words, exploiting fine-grained unstructured sparsity is possible, but it is
contingent on both the dataflow and the optimization objective, and it does not
reduce the dominant weight-transport energy.

% Table 1: structured channel harvest (#1). Complete.
\begin{table}[t]
  \centering
  \small
  \caption{Structured channel harvesting of the pruned checkpoints. The compact
    model is built by deleting emergently-dead channels and is \emph{exactly
    lossless} (maximum double-precision output deviation $<10^{-12}$; no
    fine-tuning). Throughput is measured on an NVIDIA RTX~4080 at batch $256$.
    Upper block: high-sparsity feed-forward networks, where harvesting is large.
    Lower block: networks where it is negligible---low-sparsity feed-forward
    models (below the emergence threshold) and all residual/dense models (whose
    skip/concatenation coupling leaves $0$ channels losslessly removable).}
  \label{tab:harvest}
  \begin{tabular}{l r r r r c}
    \toprule
    Model & Sparsity & \multicolumn{2}{c}{Params (M)} & Comp. & GPU $\times$ \\
    \cmidrule(lr){3-4}
          & (\%)     & full & compact & & (fp16/fp32) \\
    \midrule
    VGG-19 (CIFAR-100)   & 98.9 & 20.08 & 1.66  & $12.1\times$ & $3.1/4.2$ \\
    VGG-19 (CIFAR-10)    & 99.5 & 20.04 & 1.87  & $10.7\times$ & $2.1/2.8$ \\
    LeNet-5              & 99.5 & 0.43  & 0.046 & $9.4\times$  & $\sim\!1$\,$^{\dagger}$ \\
    VGG-19 (TinyImageNet)& 97.9 & 20.13 & 5.09  & $4.0\times$  & $2.0/2.3$ \\
    LeNet-300            & 97.9 & 0.27  & 0.097 & $2.8\times$  & $\sim\!1$\,$^{\dagger}$ \\
    \midrule
    VGG-19 (CIFAR-10)    & 90.9 & 20.04 & 18.20 & $1.1\times$  & $\sim\!1$ \\
    VGG-19 (CIFAR-100)   & 89.4 & 20.08 & 20.06 & $1.0\times$  & $\sim\!1$ \\
    VGG-19 (TinyImageNet)& 86.8 & 20.13 & 20.13 & $1.0\times$  & $\sim\!1$ \\
    ResNet-32 (CIFAR-10) & 87.9 & 1.86  & 1.86  & $1.0\times$  & $\sim\!1$ \\
    ResNet-32 (CIFAR-100)& 86.3 & 1.87  & 1.87  & $1.0\times$  & $\sim\!1$ \\
    ResNet-32 (TinyImageNet)& 89.4 & 1.89 & 1.89 & $1.0\times$ & $\sim\!1$ \\
    DenseNet-40 (CIFAR-10)& 84.5 & 1.06 & 1.06  & $1.0\times$  & $\sim\!1$ \\
    \bottomrule
  \end{tabular}

  \vspace{2pt}
  {\footnotesize $^{\dagger}$LeNet models are too small to be compute-bound on a
  GPU (kernel-launch dominated); their harvest manifests in
  parameters/energy and on CPU/edge targets rather than in GPU throughput.}
\end{table}
     % Table~\ref{tab:harvest}
% Table 2: sparsity-aware accelerator (#3, Sparseloop). SKELETON.
% VGG-19 CIFAR-100 98% is measured; the remaining rows (---) are pending the
% overnight sweep (fill from sparse_sweep_results.txt). "Energy ratio" is
% sparse/dense at the energy-optimal mapping (<1 = a reduction); "Throughput"
% is the delay-optimal dense/sparse cycle ratio (>1 = a speed-up).
\begin{table}[t]
  \centering
  \small
  \caption{Sparsity-aware accelerator (Sparseloop, $256$-PE weight-stationary
    GF4 array, $45$\,nm). Each layer is modeled as a GEMM with zero-weight
    \emph{skipping}. Feed-forward rows model the residual sparsity remaining inside
    the harvested compact network; residual/dense rows model the original sparse
    network (which cannot be harvested). Energy ratio is sparse/dense at the
    energy-optimal mapping ($<1$ is a reduction); throughput is the delay-optimal
    dense-to-sparse cycle ratio.}
  \label{tab:sparseloop}
  \begin{tabular}{l c c}
    \toprule
    Model & Energy ratio (sparse/dense) & Throughput $\times$ \\
    \midrule
    VGG-19 (CIFAR-100, 98\%)  & $0.79$ & $3.68$ \\
    VGG-19 (CIFAR-10, 99\%)   & ---    & ---    \\
    VGG-19 (TinyImageNet)     & ---    & ---    \\
    LeNet-5                   & ---    & ---    \\
    LeNet-300                 & ---    & ---    \\
    \midrule
    ResNet-32 (CIFAR-10, 95\%)   & ---    & ---    \\
    ResNet-32 (CIFAR-100, 85\%)  & ---    & ---    \\
    ResNet-32 (CIFAR-10, 86\%)   & ---    & ---    \\
    ResNet-32 (CIFAR-100, 86\%)  & ---    & ---    \\
    ResNet-32 (TinyImageNet)     & ---    & ---    \\
    DenseNet-40 (CIFAR-10)       & ---    & ---    \\
    \bottomrule
  \end{tabular}
\end{table}
  % Table~\ref{tab:sparseloop}

\paragraph{The limits of sparsity, and the case for quantization.}
Taken together, these results delineate what sparsity can and cannot do. Structured
harvesting is lossless, portable, and large, but it is available only to
feed-forward architectures above a high sparsity threshold. Accelerator-level
skipping extends the benefit to unstructured sparsity and to residual networks, but
only as a latency win on a specialized dataflow, with little or no energy
improvement because the metadata and weight-transport costs are untouched. Sparsity,
in short, is powerful but \emph{conditional}: its payoff depends on the network
topology, the sparsity level, and the target hardware.

This motivates the complementary compression axis studied in the next chapter.
Reducing the numerical \emph{precision} of the weights and activations is
\emph{architecture-agnostic}: every layer of every network---including the residual
and dense models that resist structured harvesting---shrinks by the same bit-width
ratio, with no dependence on where the nonzeros lie or how the layers are wired.
Unlike sparsity, low-precision arithmetic reduces the cost of \emph{every} operand
transfer and multiply directly, attacking exactly the weight-transport energy that
zero-skipping leaves in place. And because the two axes are orthogonal---one removes
\emph{which} values are computed, the other reduces the \emph{cost} of each value---
they compose multiplicatively: a channel-harvested, then quantized, model realizes
both the structural and the numerical savings at once. The remainder of this thesis
therefore turns from \emph{which} weights to keep to \emph{how precisely} to
represent them.
